\section{Results and Discussion}

This chapter presents the outcomes of the machine learning experiments, the final model 
performance on the expanded dataset, and evidence that the BirdSight web platform implements 
the intended citizen-science workflow. It then interprets these outcomes in relation to the 
research questions stated in the Introduction and discusses limitations and threats to validity.

\subsection{Experiment and Validation}

This part documents \emph{what was validated} and \emph{how}, before the detailed numerical and 
interface outcomes are presented in the \textit{Results} part of this chapter. The validation has 
two parts: (i) a controlled machine-learning protocol with fixed data splits and shared evaluation 
code, and (ii) a practical software check that the integrated services run together as a working 
system. The following subsubsections address each part in turn.

\subsubsection{Machine learning protocol}

The bird species classifier was developed and evaluated using a fixed experimental protocol 
to keep comparisons fair. Images were sourced from GBIF-mediated observation records with 
species-level labels, restricted to a Turkey-focused species list as described in the Methodology 
\cite{gbif2025,ebird2025}. After preprocessing and cleaning, the dataset was split into training, 
validation, and test partitions using a fixed ratio so that the test split was never used for 
parameter updates or model selection decisions reported as final performance.

Five convolutional architectures (VGG16, ResNet50, DenseNet121, MobileNetV2, InceptionV3) were
trained with ImageNet initialisation and the same high-level training recipe (two-phase transfer 
learning, shared optimiser family, and consistent evaluation code). This constitutes the primary 
controlled experiment for architecture selection. InceptionV3 was then selected and improved through 
three targeted ablation-style experiments on the original 23-class setup: label smoothing, stronger 
augmentation, and tuning the number of frozen backbone epochs. Finally, the best combined configuration 
was retrained on the expanded 30-class dataset.

Validation during training relied on held-out validation accuracy and loss, with early stopping in 
the fine-tuning phase to reduce overfitting. Final reported metrics were computed on the untouched 
test split using the same metric definitions throughout (top-1 accuracy, top-$k$ accuracy, precision
/recall/F1 summaries, and additional ranking scores as listed in 
Table~\ref{tab:model_evaluation_results} and Table~\ref{tab:final_optimized_results}).

\subsubsection{Software validation (functional checks)}

Beyond offline metrics, the integrated system was validated as a runnable end-to-end application: users 
can authenticate, submit observations with images and coordinates, request species suggestions, browse 
observations and taxonomy pages, and contribute identifications that update observation quality 
metadata. These checks confirm that the architectural split (Next.js client, Spring Boot API, FastAPI 
inference, PostgreSQL/PostGIS, and MinIO object storage) works together in a local deployment 
configuration. This validation is functional rather than a formal user study; no human-subjects 
claims are made beyond standard software engineering acceptance testing.

\subsection{Results}

This part summarises the measurable outcomes of the project. It first reports the comparative benchmark 
among five CNN architectures, then the targeted improvements applied to InceptionV3, the final performance 
after expanding to 30 species, and finally qualitative evidence from the implemented web interface (with 
references to the figures and tables already presented in the Methodology). The intent is to separate 
\emph{what was observed} from the interpretation given later in the \textit{Discussion} part of this chapter.

\subsubsection{Benchmark comparison (five architectures)}

Table~\ref{tab:model_evaluation_results} summarises the benchmark performance on the held-out test 
data. InceptionV3 achieved the strongest overall top-1 accuracy and the strongest top-3 and top-5 
behaviour among the five candidates, which is important for an interactive workflow where users pick 
among suggested species rather than only accepting the first label. DenseNet121 and ResNet50 performed 
similarly to each other and slightly below InceptionV3, while MobileNetV2 remained competitive given 
its lighter design. VGG16 lagged behind, which is consistent with the comparative literature on older 
deep stacks versus residual and factorised architectures \cite{bilgin2022bird,tang2025comparative}.

Figure~\ref{fig:model_comparison_chart} provides a compact visual comparison across models. 
Figure~\ref{fig:inception_confusion_matrix} shows that most mass lies on the diagonal for InceptionV3, 
while Figure~\ref{fig:top_confused_pairs} highlights a small set of frequently confused species pairs. 
Qualitatively, many of these pairs correspond to visually similar taxa, which matches expectations for 
fine-grained classification \cite{chen2024fetfgvc}.

\subsubsection{InceptionV3 improvement experiments}

The ablation sequence confirmed that the strongest single gain came from enhanced augmentation, with 
training and validation accuracy curves becoming more aligned compared with the baseline augmentation 
setup (Figures~\ref{fig:baseline_training_history} and~\ref{fig:enhanced_aug_training_history}). 
Freeze-epoch tuning showed a moderate effect, with 11 frozen epochs achieving the best observed 
validation accuracy in the sweep (Table~\ref{tab:freeze_epoch_experiment} and 
Figure~\ref{fig:freeze_epoch_comparison}). Label smoothing produced smaller movement in top-1 accuracy 
but supports better-calibrated probabilities, which can matter when presenting ranked suggestions to 
non-experts.

\subsubsection{Final model on the 30-class dataset}

After merging the best settings, the final InceptionV3 model was trained on 30 species and 150{,}000 
images. Table~\ref{tab:final_optimized_results} reports the resulting test performance. Despite the 
increased number of classes, the model retained strong top-1 accuracy and high top-3 accuracy, 
indicating that the ranked suggestion interface remains useful when exact top-1 classification is 
difficult.

\subsubsection{Platform and user-interface outcomes}

The implemented client covers the intended citizen-science journey. The landing experience introduces 
the workflow and citizen-science context (Figures~\ref{fig:hero_section}--\ref{fig:citizen_science}), 
while Figure~\ref{fig:user_workflow} summarises the end-to-end observation submission path.

The observation feed supports browsing at scale with filtering and pagination 
(Figure~\ref{fig:observation_feed}). The map view supports geographic exploration with clustering 
behaviour at wider zoom levels (Figure~\ref{fig:map_page}). The taxonomy browser exposes hierarchical 
navigation and taxon-scoped information (Figure~\ref{fig:taxonomy_page}).

Observation detail pages combine media, metadata, community identifications, and discussion 
(Figure~\ref{fig:observation_page}). The submission flow supports multiple images and shows model 
suggestions through a dedicated pop-up (Figures~\ref{fig:observation_submission} 
and~\ref{fig:ml_suggestions}). Together, these screenshots provide qualitative evidence that the 
platform connects model assistance with community validation rather than replacing human judgement 
entirely \cite{baker2021verification,fraisl2022citizen}.

\subsection{Discussion}

This part moves from reporting results to explaining their meaning. It connects the empirical findings 
to the research questions posed in the Introduction, relates the design choices to prior work, and 
states limitations openly (including engineering and dataset threats to validity). A short closing 
paragraph summarises the takeaway message. The discussion is kept separate from the \textit{Results} 
part so that facts and interpretations are not mixed in a way that confuses readers.

\subsubsection{Answering the research questions}

\textbf{RQ1} asked how a learning-based classifier integrated into a citizen-science platform can 
support large-scale monitoring. The results support a practical answer: automated top-$k$ suggestions 
lower the friction for labelling noisy amateur images, while the observation feed, map, and 
identification workflow collect structured, georeferenced records at scale \cite{fraisl2022citizen}. 
The quality-grade mechanism further operationalises a lightweight community consensus layer on top 
of model uncertainty, which aligns with discussions on data quality in volunteer-generated biodiversity 
datasets \cite{johnston2022outstanding}.

\textbf{RQ2} asked whether publicly available bird imagery is suitable for training deep models 
in this setting. The GBIF-mediated pipeline produced a large labelled corpus, and the final model 
achieved strong held-out performance, supporting suitability after preprocessing, splitting, and 
careful training \cite{gbif2025}. However, suitability is not ``noise-free labels'': citizen-science 
imagery varies in pose, distance, blur, and background clutter, and some species pairs remain 
partially confusable even when metrics are strong (Figure~\ref{fig:top_confused_pairs}). This is 
consistent with fine-grained recognition challenges emphasised in the literature \cite{chen2024fetfgvc}.

\subsubsection{Relation to prior work and design choices}

The benchmark ranking aligns with comparative CNN studies that frequently favour strong modern 
architectures over VGG-style stacks for challenging visual classification \cite{bilgin2022bird,
tang2025comparative,lolge2025evolution}. Separating inference into a dedicated FastAPI service matches 
common deployment guidance for keeping model lifecycles independent from core application services 
\cite{karanam2026microservice}.

\subsubsection{Limitations and threats to validity}

The scope is intentionally bounded: 30 species common in Turkey, not a global avifauna model. 
Performance metrics are computed on a specific train/val/test split; although the split is stratified 
and seeded for reproducibility, geographic or temporal leakage in aggregated public datasets can 
never be ruled out completely without additional auditing. External reverse geocoding depends on 
third-party availability and fair-use constraints \cite{nominatim2025}. Finally, the reported UI 
validation is functional rather than a controlled usability experiment with statistical participant 
outcomes.

\subsubsection{Concluding remarks}

Overall, the quantitative results demonstrate that a comparatively strong bird classifier can be 
trained from large-scale public imagery and integrated into a modular web stack. The qualitative 
screenshots demonstrate that the platform channels model output into a workflow where volunteers 
can still discuss, correct, and confirm identifications, which is central to trustworthy citizen 
science at scale \cite{baker2021verification}.

\clearpage
